\documentclass[11pt,a4paper]{article}

\usepackage[utf8]{inputenc}
\usepackage[T1]{fontenc}
\usepackage[english]{babel}
\usepackage{lmodern}
\usepackage{microtype}
\usepackage{amsmath,amssymb,mathtools}
\usepackage{geometry}
\usepackage{xcolor}
\usepackage{enumitem}
\usepackage{booktabs}
\usepackage{fancyhdr}
\usepackage[hidelinks]{hyperref}

\geometry{left=2.3cm,right=2.3cm,top=2.2cm,bottom=2.25cm,headheight=14pt}
\setlength{\parindent}{0pt}
\setlength{\parskip}{0.42em}
\setlist[itemize]{topsep=0.25em,itemsep=0.18em,leftmargin=1.5em}
\setlist[enumerate]{topsep=0.3em,itemsep=0.25em,leftmargin=1.9em}
\allowdisplaybreaks

\definecolor{azul}{RGB}{34,73,110}
\definecolor{gris}{RGB}{90,96,102}
\definecolor{azulclaro}{RGB}{238,244,249}

\pagestyle{fancy}
\fancyhf{}
\fancyhead[L]{\small\color{gris}GC2D: \texttt{ABBA\_semiimplicit}}
\fancyhead[R]{\small\color{gris}Trajectory and tangent dynamics}
\fancyfoot[C]{\thepage}
\renewcommand{\headrulewidth}{0.3pt}

\newcommand{\R}{\mathbb{R}}
\newcommand{\Id}{I}
\newcommand{\PhiAB}{\Phi^{\mathrm{ABBA}}_{h,t_n}}
\newcommand{\norm}[1]{\left\lVert #1\right\rVert}

\begin{document}

\begin{center}
  {\LARGE\bfseries\color{azul}The \texttt{ABBA\_semiimplicit} method}\par
  \vspace{0.35em}
  {\large Implicit physical step and explicit application of its Jacobian}\par
  \vspace{0.45em}
  {\small Adaptation to the GC2D problem with Hairer's symmetric projection}
\end{center}

\vspace{0.5em}

\section{Purpose of the method}

The \texttt{ABBA\_semiimplicit} method advances two objects simultaneously:
\[
 z_n\in\R^2,
 \qquad
 \xi_n\in\R^2,
\]
where $z_n$ is the physical state and $\xi_n$ is an infinitesimal perturbation or tangent
vector. The state is computed using the complete implicit ABBA method,
\[
 z_{n+1}=\Psi_{h,t_n}(z_n),
\]
whereas the tangent vector is updated by applying the exact Jacobian of that same step:
\[
 \boxed{
 \xi_{n+1}=D\Psi_{h,t_n}(z_n)\,\xi_n.
 }
\]

The term \emph{semi-implicit} is used here in a precise sense:
\begin{itemize}
 \item the trajectory $z_n\mapsto z_{n+1}$ requires solving for the implicit multiplier
       $\mu_n$;
 \item once $\mu_n$ has converged and the ABBA stages are known, the tangent dynamics is
       updated by means of a $2\times2$ linear system.
\end{itemize}

\subsection*{Essential distinction}

The Jacobian formula does not allow the physical step to be replaced by
\[
 z_{n+1}=D\Psi_{h,t_n}(z_n)z_n.
\]
This equality is false for a general nonlinear map. The correct application is
\[
 \delta z_{n+1}=D\Psi_{h,t_n}(z_n)\,\delta z_n.
\]
Therefore, $z_{n+1}$ is first computed with the implicit method, and its Jacobian is then
applied to $\xi_n$.

\section{GC2D system and Jacobian of the vector field}

The physical state is $z=(X,Y)^T$ and satisfies
\[
 \dot z=f(z,t)=J_0\nabla_z\phi(z,t),
 \qquad
 J_0=
 \begin{pmatrix}
  0&-1\\
  1&0
 \end{pmatrix}.
\]
The spatial Jacobian of the vector field is
\[
 W(z,t):=D_zf(z,t)=J_0\nabla_z^2\phi(z,t)
 =
 \begin{pmatrix}
  -\phi_{XY}&-\phi_{YY}\\
  \phi_{XX}&\phi_{XY}
 \end{pmatrix}_{(z,t)}.
\]
Since the Hessian of $\phi$ is symmetric,
\[
 W(z,t)^TJ_0+J_0W(z,t)=0.
\]

\section{Implicit ABBA physical step}

Let $s=h/2$. For a known state $z_n$ and a provisional value $\mu\in\R^2$, the two
copies are initialized as
\[
 u^{(0)}=z_n+\mu,
 \qquad
 v^{(0)}=z_n-\mu.
\]
The four $A$--$B$--$B$--$A$ stages are then applied:
\begin{align}
 u^{(1)}&=u^{(0)}+s f(v^{(0)},t_n), \tag{1a}\\
 v^{(1)}&=v^{(0)}+s f(u^{(1)},t_n), \tag{1b}\\
 v^{(2)}&=v^{(1)}+s f(u^{(1)},t_{n+1}), \tag{1c}\\
 u^{(2)}&=u^{(1)}+s f(v^{(2)},t_{n+1}). \tag{1d}
\end{align}

The multiplier for the step is the root $\mu_n=\mu_h(z_n)$ of the residual
\[
 \boxed{
 \mathcal R_n(z_n,\mu)
 :=u^{(2)}(z_n,\mu)-v^{(2)}(z_n,\mu)+2\mu=0.
 }
 \tag{2}
\]
Once the implicit iteration has converged, the next physical state is
\[
 \boxed{
 z_{n+1}=u^{(2)}+\mu_n=v^{(2)}-\mu_n.
 }
 \tag{3}
\]
The two expressions in (3) coincide because $\mathcal R_n(z_n,\mu_n)=0$.

\section{Construction of the physical Jacobian}

All quantities in this section are evaluated \emph{after} the converged root $\mu_n$ has
been obtained. Define
\begin{align*}
 W_1&=W(v^{(0)},t_n),
 &W_2&=W(u^{(1)},t_n),\\
 W_3&=W(u^{(1)},t_{n+1}),
 &W_4&=W(v^{(2)},t_{n+1}),
\end{align*}
and
\[
 S:=W_2+W_3.
\]

The Jacobian of the four ABBA stages, regarded as a map on the duplicated phase space,
is written in block form as
\[
 D\PhiAB=
 \begin{pmatrix}
  A&B\\ C&D
 \end{pmatrix},
\]
where
\begin{align*}
 A&=\Id_2+s^2W_4S,
 &B&=s(W_1+W_4)+s^3W_4SW_1,\\
 C&=sS,
 &D&=\Id_2+s^2SW_1.
\end{align*}

To account for the implicit dependence $\mu_n=\mu_h(z_n)$, introduce the four
$2\times2$ matrices
\begin{align}
 \mathcal K={}&4\Id_2-s(W_1+W_2+W_3+W_4) \notag\\
 &+s^2(W_4S+SW_1)-s^3W_4SW_1, \tag{4}\\[1mm]
 \mathcal L={}&s(W_1-W_2-W_3+W_4) \notag\\
 &+s^2(W_4S-SW_1) \notag\\
 &+s^3W_4SW_1, \tag{5}\\[1mm]
 \mathcal P={}&\Id_2+s(W_1+W_4) \notag\\*
 &+s^2W_4S+s^3W_4SW_1, \tag{6}\\[1mm]
 \mathcal Q={}&2\Id_2-s(W_1+W_4) \notag\\
 &+s^2W_4S-s^3W_4SW_1. \tag{7}
\end{align}

Here,
\[
 \mathcal K=D_\mu\mathcal R_n,
 \qquad
 \mathcal L=D_z\mathcal R_n.
\]
Differentiating the identity
\[
 \mathcal R_n\bigl(z,\mu_h(z)\bigr)=0
\]
with respect to $z$ gives
\[
 \mathcal L+\mathcal K D\mu_h(z)=0,
\]
and, provided that $\mathcal K$ is invertible,
\[
 D\mu_h(z_n)=-\mathcal K^{-1}\mathcal L.
\]
Consequently, the exact Jacobian of the implicit physical step is
\[
 \boxed{
 J_n:=D\Psi_{h,t_n}(z_n)
 =\mathcal P-\mathcal Q\mathcal K^{-1}\mathcal L.
 }
 \tag{8}
\]

The matrix $\mathcal K$ in (4) is also the Jacobian of the residual used in Newton's
method. However, $\mathcal K$ is not the physical Jacobian: the latter is $J_n$ and also
contains $\mathcal P$, $\mathcal Q$, and $\mathcal L$.

\section{Applying the Jacobian without computing an inverse}

In an implementation, $\mathcal K^{-1}$ should not be formed explicitly. To construct the
entire Jacobian, solve the system with two right-hand sides:
\[
 \boxed{
 \mathcal K\mathcal T=\mathcal L,
 \qquad
 J_n=\mathcal P-\mathcal Q\mathcal T.
 }
 \tag{9}
\]
Since $\mathcal K\in\R^{2\times2}$, its factorization is inexpensive and can be reused for
the two columns of $\mathcal L$.

If only a tangent vector $\xi_n$ is to be propagated, it is not even necessary to form
$\mathcal T$ or $J_n$. From (8),
\[
 J_n\xi_n
 =\mathcal P\xi_n-\mathcal Q\mathcal K^{-1}(\mathcal L\xi_n).
\]
Therefore, compute successively
\[
 b_n=\mathcal L\xi_n,
 \qquad
 \mathcal K y_n=b_n,
\]
and finally
\[
 \boxed{
 \xi_{n+1}=\mathcal P\xi_n-\mathcal Qy_n.
 }
 \tag{10}
\]
The only implicit part of (10) is a $2\times2$ linear system; no new nonlinear equation
appears.

\subsection{Explicit formula for the \texorpdfstring{$2\times2$}{2 x 2} system}

Let
\[
 \mathcal K=
 \begin{pmatrix}
  k_{11}&k_{12}\\k_{21}&k_{22}
 \end{pmatrix},
 \qquad
 b_n=
 \begin{pmatrix}b_1\\b_2\end{pmatrix},
 \qquad
 \Delta_{\mathcal K}=k_{11}k_{22}-k_{12}k_{21}.
\]
If $\Delta_{\mathcal K}\neq0$, the solution of $\mathcal K y_n=b_n$ is
\[
 \boxed{
 y_n=
 \frac{1}{\Delta_{\mathcal K}}
 \begin{pmatrix}
  k_{22}b_1-k_{12}b_2\\
  -k_{21}b_1+k_{11}b_2
 \end{pmatrix}.
 }
 \tag{11}
\]
This expression makes it possible to apply (10) without a general-purpose matrix solver.

\section{One-step algorithm}

Given $t_n$, $z_n$, $\xi_n$, and $h$, one step of \texttt{ABBA\_semiimplicit} is as
follows:

\begin{enumerate}
 \item Solve the nonlinear equation (2) to obtain $\mu_n$. Each residual evaluation
       executes the four ABBA stages in (1).

 \item With the converged root, recompute or retain the stages
       $u^{(0)},v^{(0)},u^{(1)},v^{(1)},v^{(2)},u^{(2)}$.

 \item Compute the new physical state using (3).

 \item Evaluate $W_1,W_2,W_3,W_4$ at the four stage points and form
       $S=W_2+W_3$.

 \item Construct $\mathcal K$, $\mathcal L$, $\mathcal P$, and $\mathcal Q$ using
       (4)--(7).

 \item Compute $b_n=\mathcal L\xi_n$ and solve $\mathcal K y_n=b_n$.

 \item Apply the Jacobian through
       $\xi_{n+1}=\mathcal P\xi_n-\mathcal Qy_n$.

 \item Advance the time: $t_{n+1}=t_n+h$.
\end{enumerate}

The output of the step is the pair
\[
 \boxed{
 (z_{n+1},\xi_{n+1}).
 }
\]

\section{Propagation of the accumulated Jacobian}

To obtain the derivative of the numerical solution with respect to the complete initial
condition, propagate a matrix $M_n\in\R^{2\times2}$ instead of a single vector. Initialize
\[
 M_0=\Id_2
\]
and, after each physical step, update it according to
\[
 \boxed{
 M_{n+1}=J_nM_n.
 }
 \tag{12}
\]
Thus,
\[
 M_N=J_{N-1}J_{N-2}\cdots J_1J_0
 =\frac{\partial z_N}{\partial z_0}.
\]
The order of the products is important: the Jacobian of the most recent step always
multiplies from the left.

The matrix update can also be performed without forming $J_n$. Solve
\[
 \mathcal K Y_n=\mathcal L M_n
\]
for the two right-hand sides and set
\[
 \boxed{
 M_{n+1}=\mathcal P M_n-\mathcal QY_n.
 }
 \tag{13}
\]

\section{Symplecticity of the tangent propagation}

If the projected ABBA step is solved exactly, its physical Jacobian satisfies
\[
 J_n^TJ_0J_n=J_0.
\]
In two dimensions, this identity is equivalent to
\[
 \det J_n=1.
\]
Therefore, the application $\xi_{n+1}=J_n\xi_n$ is a symplectic linear transformation at
each step. In addition, the accumulated product (12) is also symplectic:
\[
 M_N^TJ_0M_N=J_0,
 \qquad
 \det M_N=1.
\]

In numerical computations, it is useful to monitor
\[
 \varepsilon_{\mathrm{sp},n}
 :=\norm{J_n^TJ_0J_n-J_0},
 \qquad
 \varepsilon_{\det,n}:=\lvert\det J_n-1\rvert.
\]
These quantities will not be exactly zero because of roundoff, the tolerance used to
solve (2), and errors in the derivatives of $\phi$.

\section{Interpretation in terms of two nearby trajectories}

Let $z_n^\varepsilon=z_n+\varepsilon\xi_n$, with $\varepsilon$ small. Differentiability of
the implicit map gives
\[
 \Psi_{h,t_n}(z_n^\varepsilon)
 =\Psi_{h,t_n}(z_n)
 +\varepsilon D\Psi_{h,t_n}(z_n)\xi_n
 +\mathcal O(\varepsilon^2).
\]
Since $z_{n+1}=\Psi_{h,t_n}(z_n)$ and
$\xi_{n+1}=D\Psi_{h,t_n}(z_n)\xi_n$, it follows that
\[
 z_{n+1}^\varepsilon
 =z_{n+1}+\varepsilon\xi_{n+1}+\mathcal O(\varepsilon^2).
\]
Thus, \texttt{ABBA\_semiimplicit} simultaneously computes the reference trajectory and
the linear evolution of an infinitesimal separation from it.

\section{Implementation-ready summary}

At each time step, perform the two updates
\[
 \boxed{
 \begin{aligned}
 z_{n+1}&=\Psi_{h,t_n}(z_n)
 &&\text{using ABBA and the implicit root }\mu_n,\\[1mm]
 \xi_{n+1}&=
 \left(\mathcal P-\mathcal Q\mathcal K^{-1}\mathcal L\right)\xi_n
 &&\text{using a }2\times2\text{ linear system}.
 \end{aligned}
 }
\]
The recommended form for the second line is
\[
 b_n=\mathcal L\xi_n,
 \qquad
 \mathcal K y_n=b_n,
 \qquad
 \xi_{n+1}=\mathcal P\xi_n-\mathcal Qy_n.
\]
The physical state is not obtained by multiplying $z_n$ by the Jacobian. The Jacobian is
applied to the tangent vector or, equivalently, to the accumulated Jacobian with respect
to the initial condition.

\end{document}
